% Source: https://artofproblemsolving.com/community/c3h2976156_math_puzzles
% Stable identity: aops-topic-2976156-q02
Measure distance as a fraction of the circumference and fuel by the distance it allows, so a full tank is $\dfrac12$. Call the planes $A$, $B$, $C$; $A$ will circle the globe eastward.

\emph{Outward.} $A,B,C$ leave together eastward. At $\dfrac18$ each has $\dfrac38$; $C$ gives $\dfrac18$ to each of $A,B$ (both full) and keeps $\dfrac18$, enough to fly home. At $\dfrac14$, $A$ and $B$ each have $\dfrac38$; $B$ gives $\dfrac18$ to $A$ (full) and keeps $\dfrac14$, enough to fly home. $A$ now flies from $\dfrac14$ to $\dfrac34$ on its full tank, arriving empty.

\emph{Return escort.} Meanwhile $B$ and $C$ refuel. $C$ flies west from the island (a full tank) timed to meet $A$ at the point $\dfrac34$; $C$ then has $\dfrac14$ left, gives $\dfrac18$ to $A$, and both fly east with $\dfrac18$ each to the point $\dfrac78$. $B$ leaves the island westward, full, timed to meet them at $\dfrac78$ with $\dfrac38$ left; it gives $\dfrac18$ to each, keeping $\dfrac18$. All three then fly the last $\dfrac18$ home.

Every plane always has enough fuel for its next leg, and $A$ has flown the whole circle. $\blacksquare$ (In fact $3$ is the minimum, but that requires a separate argument.)
